%%% CHANGELOG
%%% Patch 0607 (Session 146, 2026-05-27)
%%% Meta-umbrella registration of THEO-DSL-9: O(delta^2) Orientation Phase
%%% Diagram Theorem.
%%%
%%% Synthesizes into a single citable meta-theorem the triad:
%%%   - THEO-DSL-6 (dimensional-growth structural)
%%%   - THEO-DSL-7 (edge-aligned O(delta^2) closure, Patch 0605a)
%%%   - THEO-DSL-8 (vertex-aligned O(delta^2) closure, Patch 0606a)
%%%
%%% No new proof obligations or physics introduced; this is a meta-level
%%% registration that provides a single citation handle for the complete
%%% O(delta^2) orientation phase diagram at the two specialized orientations
%%% theta_V and theta_E.
%%%
%%% Records five meta-level findings (F-DSL-9-1 through F-DSL-9-5) that
%%% are visible only at the triad level and not at the level of any single
%%% constituent umbrella.
%%%
%%% Rigor level: sketch-document Layer 3 (matches constituent umbrellas).
%%% Single-artifact discipline: 23-for-23 since Patch 0585.
%%%
%%% Session 146 closure note: with Patch 0607, the session has shipped a
%%% complete vertical from O(delta^2) closure (Patches 0605, 0606) through
%%% per-orientation umbrellas (0605a, 0606a) to the meta-umbrella (this
%%% patch). Total session output: 5 patches, 5 artifacts.

\documentclass[11pt]{article}
\usepackage{amsmath,amssymb,amsthm}
\usepackage{geometry}
\geometry{margin=1in}

\newtheorem{theorem}{Theorem}
\newtheorem{corollary}{Corollary}
\newtheorem{remark}{Remark}

\title{THEO-DSL-9: $\mathcal{O}(\delta^2)$ Orientation Phase Diagram Theorem \\[2pt]
\large Meta-Umbrella Registration (THEO-DSL-6 + THEO-DSL-7 + THEO-DSL-8)}
\author{Hyperphysics Institute --- Conscious Point Physics Programme}
\date{Patch 0607 \quad Session 146 \quad 2026-05-27}

\begin{document}
\maketitle

\begin{abstract}
This document registers THEO-DSL-9 as the meta-umbrella theorem
governing the orientation phase diagram of the $D_5$-invariant current
through second order in $\delta$. The meta-statement combines three
established theorems --- THEO-DSL-6 (dimensional-growth structural),
THEO-DSL-7 (edge-aligned $\mathcal{O}(\delta^2)$ closure), and
THEO-DSL-8 (vertex-aligned $\mathcal{O}(\delta^2)$ closure) --- into a
single citable result asserting that at the two specialized orientations
$\theta_V$ (vertex-aligned) and $\theta_E$ (edge-aligned), the current
expansion is determined through $\mathcal{O}(\delta^2)$ with explicit
subspace dimensions and coefficient values. Five meta-level findings
(F-DSL-9-1 through F-DSL-9-5) record observations that emerge only at
the triad level.
\end{abstract}

\section{Theorem Statement (THEO-DSL-9)}

\begin{theorem}[$\mathcal{O}(\delta^2)$ Orientation Phase Diagram --- THEO-DSL-9]
\label{thm:THEO-DSL-9}
Let $\mathbf{j}(\theta, \delta)$ denote the $D_5$-invariant current at
orientation $\theta$, expanded as
\[
\mathbf{j}(\theta, \delta) \;=\; \mathbf{j}_1(\theta)\,\delta
\;+\; \mathbf{j}_2(\theta)\,\delta^2 \;+\; \mathcal{O}(\delta^3).
\]
Then at the two specialized orientations $\theta \in \{\theta_V, \theta_E\}$:
\begin{enumerate}
\item[\textnormal{(i)}] \textbf{Dimensional structure} (THEO-DSL-6):
\[
\dim\,\bigl[\mathrm{span}\{\mathbf{j}_n(\theta_V)\}_{n=1,2}\bigr] \;=\; 1,
\qquad
\dim\,\bigl[\mathrm{span}\{\mathbf{j}_n(\theta_E)\}_{n=1,2}\bigr] \;=\; 2.
\]
The dimensional growth $1 \to 2$ at the $\theta_V \to \theta_E$
transition holds at \emph{both} $\mathcal{O}(\delta)$ and
$\mathcal{O}(\delta^2)$ separately (F-DSL-9-1).

\item[\textnormal{(ii)}] \textbf{Vertex-aligned closure} (THEO-DSL-8,
Patch 0606a): along $\hat{\mathbf{n}}_{\text{subst}}$,
\[
\alpha_1(\theta_V) \;=\; +\tfrac{6}{\varphi^2} \;\approx\; +2.292,
\qquad
\alpha_2(\theta_V) \;=\; -\tfrac{9}{\varphi^2} \;\approx\; -3.438,
\qquad
\frac{\alpha_2}{\alpha_1}\bigg|_{\theta_V} \;=\; -\tfrac{3}{2}.
\]

\item[\textnormal{(iii)}] \textbf{Edge-aligned closure} (THEO-DSL-7,
Patch 0605a): along the non-orthogonal basis
$\{\hat{\mathbf{n}}_\rho, \hat{\mathbf{n}}_{\text{edge}}\}$,
\[
\alpha_{1,\rho}(\theta_E),\; \alpha_{1,\text{edge}}(\theta_E)
\;\in\; \mathbb{R}_{\neq 0}
\quad\text{(Patch 0600)},
\]
\[
\alpha_{2,\rho}(\theta_E) \;=\; +\tfrac{9}{2\varphi} \;\approx\; +2.781,
\qquad
\alpha_{2,\text{edge}}(\theta_E) \;=\; +9\varphi-12 \;\approx\; +2.562
\quad\text{(Patch 0605)},
\]
with both second-order coefficients positive and both second-to-first
order ratios positive (no sign flip).
\end{enumerate}
\end{theorem}

\medskip
The rigor of THEO-DSL-9 inherits from its constituents: sketch-document
Layer 3 per DG-F15A-5 default-(A). No new proof obligations are
introduced; the meta-umbrella provides a single citation handle for the
complete result at the two specialized orientations.

\section{Constituent Umbrellas and Coverage}

\subsection*{THEO-DSL-6 --- Dimensional-Growth Structural}
Establishes the orientation-dependent dimensionality of the
$D_5$-invariant subspace: 1D at $\theta_V$, 2D at $\theta_E$. This is
the structural backbone of the phase diagram.

\subsection*{THEO-DSL-7 --- Edge-Aligned $\mathcal{O}(\delta^2)$ Closure (Patch 0605a)}
Combines Patches 0600 (first-order structural) + 0605 (second-order
closure). Establishes both edge-aligned columns of the phase diagram
with explicit non-orthogonal-basis decomposition.

\subsection*{THEO-DSL-8 --- Vertex-Aligned $\mathcal{O}(\delta^2)$ Closure (Patch 0606a)}
Combines Patches 0606 (first-order structural) + 0591 (second-order
closure). Establishes both vertex-aligned columns of the phase diagram
with closed-form scalar coefficients $\pm \kappa/\varphi^2$.

\section{The Complete $\mathcal{O}(\delta^2)$ Phase Diagram}

\begin{center}
\begin{tabular}{c|c|c}
\hline
 & \textbf{Vertex-aligned} $\theta_V$ & \textbf{Edge-aligned} $\theta_E$ \\
 & (THEO-DSL-8) & (THEO-DSL-7) \\
\hline
Subspace dim. (THEO-DSL-6) & 1D & 2D \\
& $\mathrm{span}\{\hat{\mathbf{n}}_{\text{subst}}\}$
& $\mathrm{span}\{\hat{\mathbf{n}}_\rho,\hat{\mathbf{n}}_{\text{edge}}\}$ \\[4pt]
$\mathcal{O}(\delta)$ coefficient(s)
& $\alpha_1 = +6/\varphi^2$
& $(\alpha_{1,\rho},\,\alpha_{1,\text{edge}})$ \\[2pt]
$\mathcal{O}(\delta^2)$ coefficient(s)
& $\alpha_2 = -9/\varphi^2$
& $(+9/(2\varphi),\,+9\varphi-12)$ \\[4pt]
$\alpha_2/\alpha_1$ ratio sign
& $-$ (sign flip)
& $+$ (no sign flip) \\[2pt]
Closed-form ratio
& $-3/2$ ($\varphi$-independent)
& $\varphi$-dependent, both positive \\[4pt]
$B_2(O_i)$ structure
& $\equiv 0$ identically
& non-vanishing (P0605) \\[2pt]
$B_4(O_i)$ structure
& $-\varphi/2$ uniform
& four distinct values (P0605) \\
\hline
\end{tabular}
\end{center}

\section{Meta-Level Findings}

These findings record observations visible only at the triad level
\{THEO-DSL-6, THEO-DSL-7, THEO-DSL-8\} and not within any single
constituent umbrella.

\paragraph{F-DSL-9-1: Order-Stable Dimensional Growth.}
The dimensional-growth pattern $\dim_V \to \dim_E$ is independent of
expansion order through $\mathcal{O}(\delta^2)$: at both first and
second order, $\dim_V = 1$ and $\dim_E = 2$. The dimensions of the
order-by-order invariant subspaces are not constrained a priori to
agree across orders (since the symmetric tensor square of a 1D subspace
is 1D, and that of a 2D subspace can be 2D or 3D); that they do agree
at both orientations is a non-trivial $D_5$ representation-content
property of the underlying path-class structure.

\paragraph{F-DSL-9-2: Path-Class Weights are Order-and-Orientation Invariant.}
The path-class weights $\{p_i\} = \{1,\,1/2,\,-1/(2\varphi),\,-\varphi/2\}$
established in Patch 0591 §3 are invariant under both orientation
change ($\theta_V \leftrightarrow \theta_E$, F-DSL-7-2) and order
change ($\mathcal{O}(\delta) \leftrightarrow \mathcal{O}(\delta^2)$).
What varies across the phase diagram cells is exclusively the
projection functional $S_*(O_i)$ (substrate, radial, or edge), not the
combinatorial weights.

\paragraph{F-DSL-9-3: $B_2$ Vanishing as the Vertex-Aligned Marker.}
The intra-shell cross-vertex sub-class $B_2(O_i)$ vanishes identically
at vertex-aligned orientation at both orders (F-DSL-8-2) and is
non-vanishing at edge-aligned orientation at second order
(F-DSL-7-3, Patch 0605 Remark 4.1). $B_2 \equiv 0$ is therefore a
candidate \emph{vertex-aligned marker} usable for orientation
identification in higher-order or interpolated settings.

\paragraph{F-DSL-9-4: Sign-Flip Diagnostic is Order-Stable Within Orientation.}
The vertex-aligned sign-flip $(\alpha_1 > 0,\, \alpha_2 < 0)$ and the
edge-aligned no-flip (all coefficients positive at $\mathcal{O}(\delta)$
and $\mathcal{O}(\delta^2)$) are stable across the two orders examined.
This stability supports the use of the sign-flip / no-flip signature as
an orientation-phase diagnostic at higher orders, pending verification
at $\mathcal{O}(\delta^3)$ (OPEN-DSL-9-1).

\paragraph{F-DSL-9-5: Computational Completeness of the Phase Diagram at $\mathcal{O}(\delta^2)$.}
With the triad in place, every cell of the orientation $\times$ order
phase diagram at $\{\theta_V, \theta_E\} \times \{\delta, \delta^2\}$ is
explicitly closed. No further computation is required at this order at
the two specialized orientations. The remaining computational frontier
at $\mathcal{O}(\delta^2)$ is the orientation-interpolation question
$\alpha_n(\theta)$ for $\theta \in (\theta_V, \theta_E)$ (OPEN-DSL-9-2).

\section{Registration Record}

\begin{center}
\begin{tabular}{ll}
\hline
\textbf{Registry identifier:} & THEO-DSL-9 \\
\textbf{Series:} & THEO-DSL (theorem-level package registrations) \\
\textbf{Level:} & meta-umbrella (combines three constituent theorems) \\
\textbf{Constituent theorems:} & THEO-DSL-6, THEO-DSL-7, THEO-DSL-8 \\
\textbf{Constituent patches (full chain):} & 0591, 0600, 0605, 0605a, 0606, 0606a \\
\textbf{Rigor level:} & sketch-document Layer 3 (matches constituents) \\
\textbf{Governance designation:} & DG-F15A-5 default-(A) \\
\textbf{Scope:} & $\mathcal{O}(\delta^2)$ at $\theta \in \{\theta_V, \theta_E\}$ \\
\textbf{Registration session:} & 146 \\
\textbf{Registration date:} & 2026-05-27 \\
\textbf{Single-artifact discipline:} & 23-for-23 since Patch 0585 \\
\hline
\end{tabular}
\end{center}

\section{Citation Form}

When citing this meta-package from external papers, use:

\medskip
\noindent\texttt{[THEO-DSL-9]} \emph{$\mathcal{O}(\delta^2)$ Orientation
Phase Diagram Theorem.} CPP Programme internal registry, Session 146,
2026-05-27. Constituent theorems: THEO-DSL-6, THEO-DSL-7 (Patch 0605a),
THEO-DSL-8 (Patch 0606a).

\medskip
For citation of a specific cell of the phase diagram, prefer the
constituent umbrella (THEO-DSL-7 or THEO-DSL-8) directly. Cite
THEO-DSL-9 when the cross-orientation or cross-order structure is the
subject.

\section{Forward Items}

\begin{itemize}
\item \textbf{OPEN-DSL-9-1}: Extension of the orientation phase diagram
to $\mathcal{O}(\delta^3)$. Subsumes OPEN-DSL-7-2 and OPEN-DSL-8-2.
Tests order-stability of the dimensional-growth pattern and the
sign-flip diagnostic beyond $\mathcal{O}(\delta^2)$.

\item \textbf{OPEN-DSL-9-2}: Continuous orientation interpolation
$\alpha_n(\theta)$ for $\theta \in (\theta_V, \theta_E)$ at each
order $n = 1, 2$. Subsumes OPEN-DSL-7-3 and OPEN-DSL-8-4. The
key question: how does the 1D $\to$ 2D dimensional transition occur
as $\theta$ varies continuously?

\item \textbf{OPEN-DSL-9-3}: Layer 4 algebraic hardening of constituent
orbital decompositions. Subsumes OPEN-A5-1 (Patch 0605), OPEN-A1V-1
(Patch 0606), OPEN-DSL-7-1, and OPEN-DSL-8-1.

\item \textbf{OPEN-DSL-9-4}: Identification of the physical observable
corresponding to the $\mathcal{O}(\delta^2)$ orientation phase
structure. The phase diagram is currently a mathematical object;
connecting it to a measurable in the CPP framework is the
phenomenological closure of this program.
\end{itemize}

\section{Session 146 Closure Note}

With Patch 0607, Session 146 has shipped a complete vertical on the
$\mathcal{O}(\delta^2)$ orientation phase diagram:
\begin{itemize}
\item Patch 0605 (A5\_E): edge-aligned second-order closure.
\item Patch 0605a: THEO-DSL-7 edge-aligned umbrella.
\item Patch 0606 (A1\_V): vertex-aligned first-order structural.
\item Patch 0606a: THEO-DSL-8 vertex-aligned umbrella.
\item Patch 0607 (this work): THEO-DSL-9 meta-umbrella.
\end{itemize}
The single-artifact discipline counter now stands at 23-for-23
since Patch 0585.

\end{document}
